% ============================================================
%  الحرب العالمية الثانية – دراسة تاريخية شاملة
%  مادة: الاجتماعيات – السنة الأولى باكالوريا
%  يُجمَّع باستخدام XeLaTeX
% ============================================================
\documentclass[12pt, a4paper]{article}

% ── الحزم الأساسية ──────────────────────────────────────────
\usepackage{fontspec}
\usepackage{polyglossia}
\usepackage{geometry}
\usepackage{titlesec}
\usepackage{enumitem}
\usepackage{xcolor}
\usepackage{mdframed}
\usepackage{booktabs}
\usepackage{array}
\usepackage{fancyhdr}
\usepackage{hyperref}
\usepackage{tikz}
\usepackage{microtype}

% ── إعداد اللغة ──────────────────────────────────────────────
\setmainlanguage[calendar=gregorian]{arabic}
\setotherlanguage{french}
\setmainfont{Amiri}[Script=Arabic, Scale=1.1]
\newfontfamily\arabicfont{Amiri}[Script=Arabic]

% ── الهوامش ──────────────────────────────────────────────────
\geometry{
  top=2.5cm, bottom=2.5cm,
  left=2cm, right=2.5cm,
  headheight=15pt
}

% ── الألوان ──────────────────────────────────────────────────
\definecolor{mainblue}{RGB}{0, 70, 127}
\definecolor{accentred}{RGB}{180, 30, 30}
\definecolor{lightgray}{RGB}{245, 245, 245}
\definecolor{darkgray}{RGB}{60, 60, 60}
\definecolor{goldyellow}{RGB}{200, 150, 0}

% ── تنسيق العناوين ───────────────────────────────────────────
\titleformat{\section}
  {\large\bfseries\color{mainblue}}
  {\thesection.}{0.6em}{}
  [\color{mainblue}\titlerule]

\titleformat{\subsection}
  {\normalsize\bfseries\color{accentred}}
  {\thesubsection.}{0.5em}{}

\titleformat{\subsubsection}
  {\normalsize\bfseries\color{darkgray}}
  {\thesubsubsection.}{0.5em}{}

% ── رأس وتذييل الصفحة ────────────────────────────────────────
\pagestyle{fancy}
\fancyhf{}
\fancyhead[R]{\small\color{mainblue}\textbf{الحرب العالمية الثانية – الأسباب والنتائج}}
\fancyhead[L]{\small\color{mainblue}السنة الأولى باكالوريا}
\fancyfoot[C]{\small\thepage}
\renewcommand{\headrulewidth}{0.4pt}

% ── صناديق ملونة ─────────────────────────────────────────────
\newmdenv[
  linecolor=mainblue,
  backgroundcolor=lightgray,
  linewidth=2pt,
  topline=false, bottomline=false,
  rightline=false,
  innerleftmargin=10pt,
  innerrightmargin=10pt,
  innertopmargin=8pt,
  innerbottommargin=8pt
]{sidebarbox}

\newmdenv[
  linecolor=accentred,
  backgroundcolor=lightgray,
  linewidth=1pt,
  roundcorner=4pt,
  innerleftmargin=10pt,
  innerrightmargin=10pt,
  innertopmargin=8pt,
  innerbottommargin=8pt
]{keybox}

% ────────────────────────────────────────────────────────────
\begin{document}
\thispagestyle{empty}

% ── صفحة العنوان ─────────────────────────────────────────────
\begin{center}
  {\color{mainblue}\rule{\linewidth}{2pt}}\\[6pt]
  {\LARGE\bfseries\color{mainblue} الحرب العالمية الثانية}\\[4pt]
  {\large\color{accentred} الأسباب – المراحل – النتائج}\\[6pt]
  {\color{mainblue}\rule{\linewidth}{2pt}}\\[10pt]
  {\normalsize\color{darkgray}
    مادة الاجتماعيات \quad|\quad السنة الأولى باكالوريا \quad|\quad الدورة الثانية}\\[4pt]
  {\small\color{darkgray} الفترة الزمنية: 1939 – 1945م}
\end{center}

\vspace{0.5cm}

% ── تمهيد ────────────────────────────────────────────────────
\begin{sidebarbox}
\textbf{التساؤلات المحورية:}
\begin{itemize}[leftmargin=*, nosep]
  \item ما العوامل التي أفضت إلى اندلاع الحرب العالمية الثانية؟
  \item كيف تطورت مسارات هذه الحرب عبر مرحلتيها الكبيرتين؟
  \item ما الآثار الجوهرية التي أعادت رسم معالم العالم بعد انتهائها؟
\end{itemize}
\end{sidebarbox}

\vspace{0.3cm}

\noindent
شكّل الهجوم الألماني المباغت على الأراضي البولندية في مطلع شتنبر 1939م الشرارةَ
الأولى لنزاع مسلح عالمي غير مسبوق في حجمه وتداعياته.
استقطب الصراعُ قوى كبرى من شتى القارات، ولم تُسدل ستائره إلا بعد إلقاء أول قنبلتين
نوويتين في تاريخ البشرية على المدينتين اليابانيتين هيروشيما وناغازاكي.

% ════════════════════════════════════════════════════════════
\section{أسباب الحرب العالمية الثانية}
% ════════════════════════════════════════════════════════════

\subsection{الأسباب البعيدة (غير المباشرة)}

\subsubsection{الأزمة الاقتصادية الكبرى}
تركت أزمة الكساد الكبير (1929م) وتداعياتها الاقتصادية آثاراً عميقة في بنية الاقتصادات
الأوروبية، وكان النصيب الأوفر من الضرر من حصة ألمانيا التي عانت من بطالة واسعة
وانهيار في قيمة عملتها الوطنية.

\subsubsection{شروط معاهدة فرساي 1919م}
فرضت القوى المنتصرة في الحرب الأولى على ألمانيا شروطاً مجحفة تضمنت:
\begin{itemize}[leftmargin=*, itemsep=2pt]
  \item تجريدها من مستعمراتها ومن مناطق استراتيجية داخل أوروبا.
  \item تحديد أعداد قواتها المسلحة وتقييد برامجها التسليحية.
  \item إلزامها بتعويضات مالية باهظة تحت مسمى "تعويضات الحرب".
\end{itemize}
أوجدت هذه الشروط قناعة شعبية ألمانية راسخة بضرورة مراجعتها، فتحولت إلى وقود
أذكى الحركات القومية المتطرفة.

\subsubsection{صعود الأنظمة الشمولية}
أفرزت الأزمة الاقتصادية وخيبات معاهدة فرساي بيئةً خصبة لصعود الديكتاتوريات:
\begin{itemize}[leftmargin=*, itemsep=2pt]
  \item \textbf{ألمانيا:} تولّى أدولف هتلر قيادة الحزب النازي وبلغ السلطة عام 1933م،
        فشرع في انتهاك بنود معاهدة فرساي، وضم إقليم السار (1936م)،
        والنمسا (1938م – حادثة الأنشلوس)، وإقليم السوديت بتشيكوسلوفاكيا.
  \item \textbf{إيطاليا:} رسّخ موسوليني نظامه الفاشي وسعى إلى توسيع نفوذ إيطاليا
        في البحر الأبيض المتوسط وأفريقيا.
  \item \textbf{اليابان:} اعتمد العسكريون سياسة توسعية نحو جنوب شرق آسيا
        واحتلوا مناطق واسعة من الصين.
\end{itemize}

\subsubsection{تشكّل محور برلين–روما–طوكيو}
أسفر التقارب الإيديولوجي بين الدول الثلاث عن إبرام تحالفات متعاقبة بدءاً من عام
1937م، مكّنتها من تنسيق طموحاتها التوسعية ومواجهة القوى الديمقراطية الغربية.
علاوة على ذلك، أقدمت كل من ألمانيا وإيطاليا واليابان على الانسحاب من عصبة الأمم،
مما أفقد هذه الهيئة الدولية زخمها وقدرتها على ضبط النظام الدولي.

\subsection{الأسباب المباشرة للحرب}

\subsubsection{التصعيد الألماني ومطالب هتلر الإقليمية}
أعلن هتلر مطالبته باسترجاع منطقة راينلاند وممر دانتزيغ الذي يصل ألمانيا بأراضيها
في بروسيا الشرقية، مستنداً إلى حجة الغالبية السكانية الناطقة بالألمانية.
غير أن الهدف الحقيقي تمثّل في بسط السيطرة على بولندا لضمان الموارد الزراعية
وتحقيق ما عُرف بـ"المجال الحيوي" (Lebensraum) المتجه نحو الشرق.

\subsubsection{الأحداث المتسارعة نحو الحرب}
\begin{itemize}[leftmargin=*, itemsep=3pt]
  \item نقض هتلر معاهدةَ عدم الاعتداء الألمانية-البولندية المُبرمة عام 1934م.
  \item إبرام "ميثاق الفولاذ" مع إيطاليا في مايو 1939م.
  \item توقيع الاتحاد السوفياتي وألمانيا ميثاقَ عدم الاعتداء (ميثاق مولوتوف-ريبنتروب)
        في أغسطس 1939م، الذي أزاح آخر العوائق أمام الاجتياح.
  \item اجتياح الأراضي البولندية في الأول من شتنبر 1939م، فأعلنت فرنسا وبريطانيا
        الحرب على ألمانيا في الثالث من الشهر ذاته.
\end{itemize}

% ════════════════════════════════════════════════════════════
\section{المراحل الكبرى للحرب (1939 – 1945م)}
% ════════════════════════════════════════════════════════════

\subsection{المرحلة الأولى: انتصارات دول المحور (1939 – 1942م)}

اتكأت ألمانيا في هذه المرحلة على نظرية الحرب الخاطفة (Blitzkrieg)، الجامعة بين
ضربات جوية متواصلة وزحف مدرّع سريع، مما أربك دفاعات الخصم قبل تنظيمها.

\begin{keybox}
\textbf{أبرز الأحداث والجبهات:}
\begin{itemize}[leftmargin=*, itemsep=3pt]
  \item \textbf{الجبهة الغربية (1940م):}
        اجتاحت ألمانيا بلجيكا وهولندا وفرنسا في أسابيع معدودة،
        واضطرت فرنسا إلى توقيع هدنة مع دول المحور.
  \item \textbf{معركة بريطانيا (1940م):}
        شنّت ألمانيا حملة قصف جوي مكثف على المدن البريطانية امتدت ثلاثة أشهر،
        إلا أن صمود سلاح الجو الملكي البريطاني أحبط مخطط الغزو البري.
  \item \textbf{الجبهة المتوسطية وشمال أفريقيا:}
        تقدمت القوات الإيطالية الألمانية نحو مصر، في حين سيطرت إيطاليا على
        ألبانيا ويوغوسلافيا واليونان.
  \item \textbf{الجبهة السوفيتية (1941م):}
        أطلقت ألمانيا عملية بارباروسا بغزو مفاجئ للاتحاد السوفياتي، ودمّرت
        معظم سلاحه الجوي على أرض المدارج، وتقدمت حتى ضواحي موسكو قبل أن
        يوقفها الشتاء القارس والمقاومة السوفيتية العنيدة.
  \item \textbf{المحيط الهادئ (ديسمبر 1941م):}
        نفّذت اليابان ضربة استباقية ضد القاعدة البحرية الأمريكية في بيرل هاربر،
        مما دفع الولايات المتحدة إلى الانخراط الرسمي في الحرب.
\end{itemize}
\end{keybox}

\subsection{المرحلة الثانية: انقلاب الموازين لصالح الحلفاء (1942 – 1945م)}

شكّل عام 1942م نقطة تحول استراتيجية حاسمة، إذ بدأت كفة الحلفاء ترجح بفضل
التفوق اللوجستي الأمريكي والصمود السوفيتي والعمليات المنسقة على الجبهات كافة.

\begin{itemize}[leftmargin=*, itemsep=4pt]
  \item \textbf{معركة العلمين (1942م):}
        حقق الحلفاء انتصاراً حاسماً في مصر، ثم تلاقت جيوشهم القادمة
        من الدار البيضاء مع القوات المتقدمة من الشرق، فطُرد المحور
        نهائياً من ليبيا ومصر وتونس عام 1943م.
  \item \textbf{معركة ستالينغراد (1943م):}
        تُعدّ من أطول المعارك البرية في التاريخ وأكثرها دموية؛ انتهت بهزيمة
        ألمانيا وفقدانها مئات الآلاف من جنودها أسرى ومقاتلين،
        وشكّلت نقطة الارتداد الألماني على الجبهة الشرقية.
  \item \textbf{الإنزال في نورماندي (يونيو 1944م):}
        فتح التحالف الكبير (الولايات المتحدة، بريطانيا، الاتحاد السوفياتي) جبهة
        غربية بإنزال ضخم، ما أفضى إلى تحرير فرنسا ثم التقدم نحو برلين.
  \item \textbf{استسلام ألمانيا (مايو 1945م):}
        وقّعت ألمانيا على وثيقة الاستسلام الرسمي في الثامن من مايو 1945م.
  \item \textbf{المسرح الآسيوي والقنبلتان الذريتان:}
        رفضت اليابان مطالب الاستسلام، فألقت الولايات المتحدة قنبلة نووية
        على هيروشيما في 6 غشت 1945م، وأخرى على ناغازاكي في 9 غشت،
        مما أجبر اليابان على الاستسلام اللاشرطي في 15 غشت 1945م.
\end{itemize}

% ════════════════════════════════════════════════════════════
\section{نتائج الحرب العالمية الثانية}
% ════════════════════════════════════════════════════════════

\subsection{الخسائر البشرية والمادية}

\subsubsection{الخسائر البشرية}
راح ضحية هذا النزاع ما يتراوح بين خمسة وسبعين ومئة مليون شخص بين مدني وعسكري،
وكانت أعلى الفواتير البشرية من نصيب الاتحاد السوفياتي وبولندا وألمانيا ويوغوسلافيا.
أفرز ذلك:
\begin{itemize}[leftmargin=*, itemsep=2pt]
  \item انحساراً ديمغرافياً حاداً وتراجعاً في نسب المواليد.
  \item اختلالاً في الهرم العمري لسنوات طويلة لاحقة.
  \item موجات نزوح وهجرة غيّرت التوزيع السكاني لدول عديدة.
\end{itemize}

\subsubsection{الخسائر المادية}
طالت يد الدمار المدن الأوروبية الكبرى وبنيتها التحتية من مصانع وطرق وجسور وشبكات:
\begin{itemize}[leftmargin=*, itemsep=2pt]
  \item تراجع حاد في مؤشرات الإنتاج الصناعي والزراعي.
  \item ارتفاع مستويات الديون الخارجية لإعادة الإعمار.
  \item تضخم اقتصادي وتدني ملحوظ في مستوى معيشة السكان.
\end{itemize}

في المقابل، خرجت الولايات المتحدة الأمريكية من الحرب أكثر اقتصاديات العالم نمواً،
إذ وفّرت لآلتها الإنتاجية طلباً متواصلاً لتلبية الاحتياجات العسكرية والغذائية
للدول المتحاربة، مما رسّخ زعامتها الاقتصادية العالمية.

\subsection{التحولات السياسية والجيوسياسية}

\subsubsection{إعادة رسم خريطة أوروبا}
\begin{itemize}[leftmargin=*, itemsep=3pt]
  \item عُقد مؤتمرا يالطا وبوتسدام عامَي 1945م للاتفاق على النظام الدولي الجديد.
  \item توسّعت مساحات الاتحاد السوفياتي وبولندا وبلغاريا وفرنسا ويوغوسلافيا،
        فيما تقلصت أراضي اليابان وألمانيا والنمسا.
  \item قُسّمت ألمانيا إلى أربع مناطق نفوذ تحت إشراف الدول الأربع:
        الولايات المتحدة، بريطانيا، فرنسا، والاتحاد السوفياتي،
        وطُبّق التقسيم ذاته على عاصمتَي برلين وفيينا.
\end{itemize}

\subsubsection{نشأة منظمة الأمم المتحدة}
حرصاً على تفادي كارثة مسلحة ثالثة، تداعت الدول الكبرى إلى مؤتمر سان فرانسيسكو
عام 1945م حيث أسّست منظمة الأمم المتحدة خلفاً لعصبة الأمم المنهكة، وأنيطت بها:
\begin{itemize}[leftmargin=*, itemsep=2pt]
  \item \textbf{الجمعية العامة:} المنتدى التداولي لجميع الدول الأعضاء.
  \item \textbf{مجلس الأمن:} المسؤول الأول عن صون السلم والأمن الدوليين.
  \item \textbf{محكمة العدل الدولية:} الجهاز القضائي لتسوية النزاعات.
  \item \textbf{وكالات متخصصة:} منظمة الصحة العالمية، منظمة العمل الدولية،
        صندوق النقد الدولي، والبنك الدولي.
\end{itemize}

\begin{sidebarbox}
\textbf{أهداف ميثاق الأمم المتحدة:}
\begin{itemize}[leftmargin=*, nosep]
  \item صون السلم الدولي ومنع الحروب بالطرق الدبلوماسية والتحكيمية.
  \item تعزيز حقوق الإنسان وضمان الحريات الأساسية دون تمييز.
  \item تشجيع التعاون الدولي في المجالات الاقتصادية والاجتماعية والثقافية.
\end{itemize}
\end{sidebarbox}

% ════════════════════════════════════════════════════════════
\section{خاتمة: عالم ما بعد الحرب وبوادر الحرب الباردة}
% ════════════════════════════════════════════════════════════

أفرزت الحرب العالمية الثانية نظاماً دولياً ثنائي القطبية:
\begin{itemize}[leftmargin=*, itemsep=3pt]
  \item \textbf{المعسكر الغربي الرأسمالي} بقيادة الولايات المتحدة الأمريكية،
        المتمسّك بمبدأ اقتصاد السوق والديمقراطية الليبرالية.
  \item \textbf{المعسكر الشرقي الاشتراكي} بزعامة الاتحاد السوفياتي،
        المعتمِد على التخطيط المركزي والأيديولوجيا الماركسية-اللينينية.
\end{itemize}

\noindent
أشعل التنافس الأيديولوجي والاستراتيجي بين القطبين فتيلَ ما عُرف بـ\textbf{"الحرب الباردة"}،
التي اتسمت بسباق تسليح محموم، وتشكيل أحلاف عسكرية متقابلة (حلف الناتو في مواجهة
حلف وارسو)، وتنافس على مناطق النفوذ في آسيا وأفريقيا وأمريكا اللاتينية،
وذلك في ظل غياب مواجهة مسلحة مباشرة بين القوتين النوويتين الكبريين.

\vspace{1cm}
\begin{center}
  {\color{mainblue}\rule{0.5\linewidth}{0.8pt}}\\[4pt]
  {\small\color{darkgray}
    انتهى الدرس – بالتوفيق في امتحاناتكم}\\[2pt]
  {\color{mainblue}\rule{0.5\linewidth}{0.8pt}}
\end{center}

\end{document}
